% !TeX root = Linear Algebra.tex

\usepackage{amsmath, amsfonts, amssymb, amsthm}
\usepackage{graphicx}
\usepackage{physics}
\usepackage{tikz}
\usepackage{mathtools}
\usepackage{xspace}

\usepackage{enumitem}
\setlist{noitemsep}

\usepackage{scrlayer-scrpage}
\ohead{\color{blue!35!black} \scshape VM}
\cfoot*{\pagemark}
\setlength{\headheight}{15pt}

\usepackage{tocloft}
\renewcommand{\cftdot}{}

\usepackage{xcolor}
\usepackage{hyperref}
\definecolor{linkcolor}{RGB}{32, 96, 192}
\hypersetup{colorlinks, linkcolor = linkcolor, urlcolor = linkcolor,
  citecolor = red, linktocpage = true}
\usepackage{bookmark}
\bookmarksetup{color = [RGB]{32, 96, 192}}

\usepackage[capitalise]{cleveref}
\crefname{Theorem}{theorem}{theorems}
\crefname{Lemma}{lemma}{lemmas}
\crefname{Corollary}{corollary}{corollaries}
\crefname{Definition}{definition}{definitions}
\crefname{Example}{example}{examples}
\crefname{Exercise}{exercise}{exercises}
\usepackage[toc, page]{appendix}

\usepackage{newpxmath}
\usepackage{charter}
\usepackage[T1]{fontenc}

\newtheorem{Theorem}{Theorem}[section]
\newtheorem{Lemma}[Theorem]{Lemma}
\newtheorem{Corollary}[Theorem]{Corollary}

\theoremstyle{definition}
\newtheorem{Definition}[Theorem]{Definition}
\newtheorem*{Definition*}{Definition}
\newtheorem{Example}[Theorem]{Example}
\newtheorem*{Example*}{Example}
\newtheorem{Exercise}{Exercise}[section]

\theoremstyle{remark}
\newtheorem*{Remark*}{Remark}
\newtheorem*{Solution*}{Solution}
\newtheorem*{Note*}{Note}

\definecolor{alertcolor}{rgb}{0.8, 0.5, 0}
\newcommand{\newterm}[1]{{\color{alertcolor} #1}}
\definecolor{notecolor}{rgb}{0, 0.5, 0.1}
\newcommand{\note}[1]{{\noindent\footnotesize\color{notecolor}\textbf{Note.} #1}}

\DeclareMathOperator{\im}{im}
\DeclareMathOperator{\nullity}{nullity}
\let\Re\relax
\DeclareMathOperator{\Re}{Re}
\let\Im\relax
\DeclareMathOperator{\Im}{Im}
\newcommand{\id}{\mathrm{id}}
\newcommand{\nth}{\textsuperscript{th}\xspace}
